\documentclass[12pt,a4paper]{article}
\usepackage[utf8]{inputenc}
\usepackage[french]{babel}
\usepackage{graphicx}
\usepackage{amsmath}
\usepackage{amssymb}
\usepackage{hyperref}
\usepackage{float}
\usepackage{caption}
\usepackage{geometry}
\usepackage{array}
\usepackage{booktabs}
\geometry{margin=2.5cm}

\hypersetup{
    colorlinks=true,
    linkcolor=blue,
    urlcolor=blue,
    citecolor=blue
}

\title{Propulsion Spatiale par Effet Hall :\\
Électronique de Puissance, Technologies et Applications}
\author{Dylan Perinetti \and Hamet Doumbouya}
\date{\today}

\begin{document}

\maketitle

\begin{abstract}
Dans quelle mesure l'électronique de puissance conditionne-t-elle les performances, la stabilité et le rendement des moteurs à effet Hall dans les systèmes de propulsion spatiale ? Ce rapport analyse les principes physiques des propulseurs à effet Hall, l'architecture de leur chaîne de conversion d'énergie, le rôle central de l'électronique de puissance (PPU), ainsi que leurs applications et perspectives de développement.
\end{abstract}

\tableofcontents
\newpage

%% ============================================================
%% I. INTRODUCTION ET CONTEXTE HISTORIQUE
%% ============================================================
\section{Introduction et contexte historique}

\subsection{De la propulsion chimique à la propulsion électrique}

La propulsion spatiale repose traditionnellement sur des moteurs chimiques, qui produisent une forte poussée mais consomment d'importantes quantités de propergol. Dès le milieu du XX\textsuperscript{e} siècle, une alternative a émergé : la propulsion électrique. Plutôt que de brûler un carburant, celle-ci utilise l'énergie électrique (fournie par des panneaux solaires ou un réacteur nucléaire) pour accélérer un propergol ionisé à très haute vitesse.

L'intérêt principal réside dans l'\textbf{impulsion spécifique} ($I_{sp}$) : un moteur chimique classique atteint 300--450~s, tandis qu'un propulseur à effet Hall peut dépasser 1500--3000~s. Cela signifie qu'à masse de propergol égale, un propulseur électrique peut fournir un $\Delta V$ (variation de vitesse) bien plus important, rendant possibles des missions longue durée avec moins de carburant embarqué.

\subsection{Naissance des moteurs à effet Hall}

Les propulseurs à effet Hall ont été développés en \textbf{Union Soviétique dans les années 1960}, principalement au sein de l'Institut Kurchatov et du centre de recherche Fakel. Le premier vol spatial d'un propulseur de type SPT (Stationary Plasma Thruster) a lieu en 1972 à bord du satellite Meteor.

Pendant des décennies, cette technologie est restée confidentielle côté occidental. Ce n'est qu'à partir des \textbf{années 1990}, après la chute de l'URSS, que les laboratoires européens et américains ont pu accéder à ces travaux. Depuis, l'adoption a été massive : les propulseurs à effet Hall équipent aujourd'hui aussi bien les satellites de télécommunications que les sondes interplanétaires.

\subsection{Lien avec l'électronique de puissance}

Le fonctionnement d'un propulseur à effet Hall dépend étroitement de l'\textbf{électronique de puissance}. En effet, le propulseur ne fonctionne pas en autonomie : il nécessite une unité de traitement de puissance (PPU, \textit{Power Processing Unit}) qui convertit l'énergie brute des panneaux solaires en tensions et courants précisément régulés pour alimenter la décharge plasma, les bobines magnétiques et la cathode neutralisatrice. La qualité de cette conversion conditionne directement les performances, la stabilité et le rendement du propulseur.

\subsection{Objectifs du document}

Ce rapport analyse le fonctionnement des propulseurs à effet Hall en suivant la chaîne complète, de la physique du plasma jusqu'à l'électronique de puissance :
\begin{itemize}
    \item Les principes physiques fondamentaux (plasma, effet Hall, accélération ionique)
    \item L'architecture système et la chaîne de conversion d'énergie
    \item L'électronique de puissance : topologies, contraintes, régulation
    \item Les performances, applications et cas d'étude
    \item Les limites, verrous technologiques et perspectives
\end{itemize}

%% ============================================================
%% II. PHYSIQUE DES PLASMAS ET PRINCIPE DU MOTEUR À EFFET HALL
%% ============================================================
\section{Physique des plasmas et principe du moteur à effet Hall}

\subsection{Qu'est-ce qu'un plasma ?}

Un plasma est un \textbf{gaz ionisé}, c'est-à-dire un état de la matière dans lequel une fraction significative des atomes a perdu un ou plusieurs électrons. Il est constitué d'un mélange d'ions positifs, d'électrons libres et d'atomes neutres. Bien que globalement quasi-neutre (autant de charges positives que négatives), le plasma possède des propriétés très particulières :

\begin{table}[H]
\centering
\caption{Propriétés fondamentales du plasma}
\begin{tabular}{>{\bfseries}l p{10cm}}
\toprule
Propriété & Description \\
\midrule
Quasi-neutralité & Le nombre de charges positives est approximativement égal au nombre de charges négatives : $n_e \approx n_i$. \\
Mobilité différentielle & Les électrons, beaucoup plus légers que les ions ($m_e \ll m_i$), sont nettement plus mobiles. \\
Comportement collectif & Les particules interagissent via les champs électromagnétiques qu'elles génèrent collectivement. \\
Conductivité électrique & Le plasma conduit le courant électrique grâce à la présence de porteurs de charges libres. \\
\bottomrule
\end{tabular}
\end{table}

Dans un propulseur à effet Hall, le propergol utilisé est généralement du \textbf{xénon} (Xe), un gaz noble choisi pour sa masse atomique élevée (131~u), sa faible énergie d'ionisation (12,1~eV) et sa non-réactivité chimique, ce qui préserve les composants du propulseur.

\subsection{L'effet Hall : principe fondamental}

L'effet Hall, découvert en 1879 par Edwin Hall, décrit le phénomène suivant : lorsqu'un courant électrique traverse un conducteur soumis à un champ magnétique perpendiculaire, une \textbf{différence de potentiel transverse} (tension de Hall) apparaît dans la direction perpendiculaire à la fois au courant et au champ magnétique.

Dans le contexte du propulseur, ce principe se manifeste de la manière suivante :

\begin{enumerate}
    \item Un \textbf{champ électrique axial} $\vec{E}$ est établi entre l'anode (au fond du canal) et la cathode (à la sortie). Ce champ accélère les ions vers la sortie.
    \item Un \textbf{champ magnétique radial} $\vec{B}$ est créé par des bobines électromagnétiques. Ce champ est perpendiculaire au champ électrique.
    \item Sous l'action croisée de $\vec{E}$ et $\vec{B}$, les \textbf{électrons} sont piégés dans un mouvement de \textbf{dérive azimutale} (rotation autour de l'axe du canal). La force de Lorentz les empêche de traverser directement le canal :
    \begin{equation}
    \vec{F} = q(\vec{E} + \vec{v} \times \vec{B})
    \end{equation}
    \item Ce mouvement de dérive des électrons constitue le \textbf{courant de Hall}. Les électrons, piégés dans cette zone, augmentent considérablement la probabilité de collision avec les atomes neutres de xénon, favorisant ainsi l'ionisation.
\end{enumerate}

\textbf{Point clé :} Les ions, beaucoup plus lourds que les électrons, ne sont quasiment pas affectés par le champ magnétique (leur rayon de Larmor est bien plus grand que les dimensions du canal). Ils sont donc accélérés librement par le champ électrique axial vers la sortie du propulseur, produisant la poussée.

\begin{figure}[H]
\centering
\includegraphics[width=0.75\textwidth]{images/hall-thruster-schema.png}
\caption{Schéma de principe d'un propulseur à effet Hall. Le champ électrique axial $\vec{E}$ accélère les ions vers la sortie, tandis que le champ magnétique radial $\vec{B}$ piège les électrons dans une dérive azimutale (courant de Hall), maximisant l'ionisation du xénon. Source~: David Staack, domaine public, via Wikimedia Commons.}
\label{fig:hall-schema}
\end{figure}

\subsection{Processus d'ionisation}

L'ionisation se produit par \textbf{collision électronique}. Les électrons piégés dans la zone de décharge possèdent une énergie cinétique suffisante (typiquement 20--30~eV) pour arracher un électron aux atomes neutres de xénon lors d'une collision :

\begin{equation}
\text{Xe} + e^- \longrightarrow \text{Xe}^+ + 2e^-
\end{equation}

Ce processus est auto-entretenu : chaque ionisation produit un électron supplémentaire qui peut à son tour ioniser d'autres atomes. Le taux d'ionisation dépend de :
\begin{itemize}
    \item La densité du gaz neutre dans la zone de décharge
    \item L'énergie des électrons (contrôlée par la tension de décharge)
    \item L'intensité du champ magnétique (qui détermine le temps de confinement des électrons)
\end{itemize}

\subsection{Accélération des ions et génération de la poussée}

Les ions Xe$^+$ créés dans la zone d'ionisation sont accélérés par le champ électrique axial. Un ion de charge $q$ traversant une différence de potentiel $V_d$ (tension de décharge) acquiert une énergie cinétique :

\begin{equation}
\frac{1}{2} m_i v_e^2 = q \cdot V_d
\end{equation}

d'où la vitesse d'éjection :

\begin{equation}
v_e = \sqrt{\frac{2 q V_d}{m_i}}
\end{equation}

Pour une tension de décharge typique de $V_d = 300$~V et du xénon ($m_i = 2{,}18 \times 10^{-25}$~kg), on obtient $v_e \approx 20$~km/s.

La \textbf{poussée} $F$ est donnée par la relation fondamentale :

\begin{equation}
F = \dot{m} \cdot v_e
\end{equation}

où $\dot{m}$ est le débit massique de propergol. L'\textbf{impulsion spécifique}, qui mesure l'efficacité du propulseur (combien de temps une unité de masse de propergol peut produire une unité de poussée), s'exprime :

\begin{equation}
I_{sp} = \frac{v_e}{g_0} = \frac{1}{g_0}\sqrt{\frac{2 q V_d}{m_i}}
\end{equation}

où $g_0 = 9{,}81$~m/s$^2$ est l'accélération gravitationnelle standard. Avec $v_e = 20$~km/s, on obtient $I_{sp} \approx 2040$~s.

Le \textbf{rendement global} du propulseur est défini comme le rapport entre la puissance cinétique du faisceau d'ions et la puissance électrique fournie :

\begin{equation}
\eta = \frac{\frac{1}{2}\dot{m} v_e^2}{P_{\text{élec}}}
\end{equation}

%% ============================================================
%% III. ARCHITECTURE SYSTÈME ET CHAÎNE DE CONVERSION D'ÉNERGIE
%% ============================================================
\section{Architecture système et chaîne de conversion d'énergie}

\subsection{Vue d'ensemble de la chaîne énergétique}

Un système de propulsion à effet Hall ne se résume pas au seul propulseur : il s'inscrit dans une \textbf{chaîne de conversion d'énergie} complète, depuis la source primaire (panneaux solaires) jusqu'à l'éjection des ions. Cette chaîne se décompose comme suit :

\begin{enumerate}
    \item \textbf{Source d'énergie} : les panneaux solaires produisent un courant continu (bus DC), typiquement à 28~V, 42~V ou 100~V selon le satellite.
    \item \textbf{PPU (Power Processing Unit)} : cette unité convertit la tension du bus en plusieurs sorties régulées adaptées aux besoins du propulseur.
    \item \textbf{Propulseur} : le moteur à effet Hall proprement dit, avec sa décharge plasma.
    \item \textbf{Systèmes auxiliaires} : réservoir de xénon, vannes de régulation du débit, cathode neutralisatrice.
\end{enumerate}

\subsection{Architectures des propulseurs}

Deux architectures principales dominent :

\begin{itemize}
    \item \textbf{SPT (Stationary Plasma Thruster)} : le canal de décharge est annulaire, avec des parois en céramique diélectrique (nitrure de bore, BN-SiO$_2$). Le plasma interagit avec les parois, ce qui contribue à l'ionisation mais provoque aussi de l'érosion. C'est l'architecture la plus répandue (SPT-100, PPS-1350).

    \item \textbf{TAL (Thruster with Anode Layer)} : le canal est plus court et les parois sont métalliques (conductrices). L'ionisation et l'accélération se concentrent dans une couche mince près de l'anode. Cette configuration réduit l'érosion mais présente des caractéristiques de décharge différentes.
\end{itemize}

\begin{figure}[H]
\centering
\includegraphics[width=0.7\textwidth]{images/spt-architecture.png}
\caption{Architecture d'un propulseur SPT : canal annulaire en céramique diélectrique, bobines magnétiques internes et externes, anode servant aussi de distributeur de gaz.}
\label{fig:spt-arch}
\end{figure}

\begin{figure}[H]
\centering
\includegraphics[width=0.7\textwidth]{images/tal-architecture.png}
\caption{Architecture d'un propulseur TAL : canal court à parois métalliques, couche d'anode mince, ionisation concentrée.}
\label{fig:tal-arch}
\end{figure}

\subsection{La PPU : interface entre source et propulseur}

La \textbf{Power Processing Unit} (PPU) est l'élément central de la chaîne. Elle doit fournir simultanément plusieurs sorties régulées :

\begin{table}[H]
\centering
\caption{Sorties de la PPU et leurs caractéristiques}
\begin{tabular}{>{\bfseries}l c c p{5cm}}
\toprule
Sous-système alimenté & Tension & Puissance & Rôle \\
\midrule
Circuit de décharge & 200--400~V & \textasciitilde~kW & Maintien du plasma et accélération des ions \\
Bobines magnétiques & 10--30~V & Dizaines de W & Création du champ magnétique radial \\
Cathode neutralisatrice & 10--15~V & \textasciitilde~50~W & Émission d'électrons pour neutraliser le faisceau \\
Allumeur de cathode & Variable & Variable & Amorçage de la cathode au démarrage \\
\bottomrule
\end{tabular}
\end{table}

La PPU doit assurer une \textbf{isolation galvanique} entre le bus satellite et les circuits du propulseur, pour protéger l'électronique embarquée des perturbations liées au plasma.

%% ============================================================
%% IV. ÉLECTRONIQUE DE PUISSANCE (CŒUR DU SUJET)
%% ============================================================
\section{Électronique de puissance}

L'électronique de puissance constitue le \textbf{cœur technologique} qui rend possible l'exploitation des propulseurs à effet Hall. Sans une PPU performante, le propulseur ne peut pas fonctionner de manière stable ni efficace.

\subsection{Topologies de conversion DC/DC}

La PPU utilise des \textbf{convertisseurs DC/DC} pour adapter la tension du bus satellite aux différents besoins du propulseur. Les principales topologies employées sont :

\begin{itemize}
    \item \textbf{Convertisseur Buck (abaisseur)} : utilisé pour alimenter les bobines magnétiques et la cathode, qui nécessitent des tensions inférieures à celle du bus.
    \item \textbf{Convertisseur Boost (élévateur)} : utilisé lorsque la tension du bus est inférieure à la tension de décharge requise (200--400~V).
    \item \textbf{Convertisseur Full-Bridge isolé} : topologie privilégiée pour le circuit de décharge principal. Elle assure l'isolation galvanique via un transformateur haute fréquence et permet de transférer des puissances de l'ordre du kW avec un bon rendement.
\end{itemize}

Les convertisseurs fonctionnent en \textbf{découpage} (commutation à haute fréquence, typiquement 50--200~kHz) avec une commande par modulation de largeur d'impulsion (PWM). Des techniques de \textbf{commutation douce} (ZVS, ZCS) sont souvent utilisées pour réduire les pertes de commutation et les perturbations électromagnétiques.

\subsection{Contraintes de conception spatiale}

La conception de l'électronique de puissance pour l'espace est soumise à des contraintes bien plus sévères que pour les applications terrestres :

\begin{table}[H]
\centering
\caption{Contraintes de conception de la PPU spatiale}
\begin{tabular}{>{\bfseries}l p{9cm}}
\toprule
Contrainte & Exigence \\
\midrule
Rendement & $> 90\%$, idéalement 93--95\%. Chaque pourcent perdu se traduit en chaleur à dissiper dans le vide. \\
Masse & Minimale : chaque kilogramme envoyé en orbite coûte plusieurs milliers d'euros. \\
Thermique & Dissipation uniquement par rayonnement (pas de convection dans le vide spatial). \\
Radiations & Composants \textit{rad-hard} (durcis aux radiations), redondance des circuits critiques. \\
CEM & Filtrage EMI rigoureux : le plasma génère des perturbations qui peuvent affecter l'électronique du satellite. \\
Fiabilité & Durée de vie de 15--20 ans sans maintenance possible. \\
\bottomrule
\end{tabular}
\end{table}

\subsection{Commande et régulation}

Le système de régulation de la PPU assure le contrôle précis des grandeurs électriques fournies au propulseur. Il comprend :

\begin{itemize}
    \item Des \textbf{boucles de rétroaction} courant/tension utilisant des régulateurs PI ou PID pour maintenir la tension de décharge et le courant à leurs valeurs de consigne.
    \item Un \textbf{contrôle adaptatif} qui ajuste les paramètres en fonction de l'évolution de l'impédance du plasma au cours de la vie du propulseur (l'érosion des parois modifie la géométrie du canal).
    \item Des \textbf{protections} contre les surtensions, surintensités et extinctions de la décharge, avec des séquences de rallumage automatique.
\end{itemize}

\subsection{Couplage plasma--électronique : un point critique}

Le plasma n'est pas une charge passive : son comportement dynamique influence directement la conception de la PPU. Trois phénomènes majeurs interviennent :

\begin{enumerate}
    \item \textbf{Oscillations de décharge} (mode ``respiration'', 10--30~kHz) : le plasma oscille périodiquement entre des phases d'ionisation intense et des phases de faible ionisation. La PPU doit filtrer ces oscillations pour ne pas les amplifier.
    \item \textbf{Impédance variable} : le plasma présente une impédance non linéaire qui varie avec les conditions de fonctionnement (débit, tension, usure). Les boucles de régulation doivent être dimensionnées avec des marges de stabilité suffisantes.
    \item \textbf{Transitoires de démarrage et d'extinction} : l'allumage et l'extinction du plasma créent des transitoires violents que la PPU doit gérer sans dommage.
\end{enumerate}

%% ============================================================
%% V. PERFORMANCES ET MODÉLISATION
%% ============================================================
\section{Performances et modélisation}

\subsection{Paramètres de performance}

\begin{table}[H]
\centering
\caption{Caractéristiques typiques des propulseurs à effet Hall}
\begin{tabular}{>{\bfseries}l c}
\toprule
Paramètre & Valeur typique \\
\midrule
Puissance électrique & 0,5 -- 20~kW \\
Poussée & 20 -- 500~mN \\
Impulsion spécifique ($I_{sp}$) & 1500 -- 3000~s \\
Vitesse d'éjection & 15 -- 30~km/s \\
Rendement global & 45 -- 60\% \\
Tension de décharge & 200 -- 400~V \\
Durée de vie & 5000 -- 10\,000~h \\
\bottomrule
\end{tabular}
\end{table}

À titre de comparaison, un moteur chimique typique produit des poussées de l'ordre du kN (1000 fois plus) mais avec une $I_{sp}$ de seulement 300~s. Le propulseur à effet Hall est donc inadapté aux phases de lancement, mais idéal pour les \textbf{manœuvres longue durée} où l'efficacité prime sur la poussée instantanée.

\subsection{Méthodes de modélisation}

La conception et l'optimisation des propulseurs font appel à plusieurs approches complémentaires :

\begin{itemize}
    \item \textbf{Modélisation analytique} : équations de conservation (masse, quantité de mouvement, énergie) appliquées au plasma. Ces modèles simplifiés permettent un dimensionnement rapide.
    \item \textbf{Simulation numérique} : les codes PIC (\textit{Particle-In-Cell}) suivent individuellement les trajectoires des particules chargées, tandis que les codes hybrides traitent les électrons comme un fluide et les ions comme des particules. Ces simulations permettent de comprendre les instabilités du plasma.
    \item \textbf{Essais expérimentaux} : les propulseurs sont testés dans des \textbf{caissons à vide} reproduisant les conditions spatiales (pression $< 10^{-5}$~mbar). Des diagnostics optiques (spectroscopie, sondes de Langmuir) mesurent les propriétés du plasma.
\end{itemize}

%% ============================================================
%% VI. APPLICATIONS ET CAS D'ÉTUDE
%% ============================================================
\section{Applications et cas d'étude}

\subsection{Satellites géostationnaires}

Les satellites géostationnaires utilisent les propulseurs à effet Hall pour le \textbf{maintien à poste} (\textit{station-keeping}) : les corrections régulières d'orbite nécessaires pour compenser les perturbations gravitationnelles (Lune, Soleil, non-sphéricité de la Terre). La propulsion électrique permet de réduire la masse de propergol de 40 à 50\% par rapport à la propulsion chimique, ce qui se traduit soit par un satellite plus léger, soit par une durée de vie allongée.

\subsection{Missions d'exploration interplanétaire}

\begin{table}[H]
\centering
\caption{Missions spatiales utilisant des propulseurs à effet Hall}
\begin{tabular}{>{\bfseries}l l l}
\toprule
Mission & Agence & Application \\
\midrule
SMART-1 & ESA & Démonstrateur technologique vers la Lune \\
BepiColombo & ESA/JAXA & Propulsion primaire vers Mercure \\
Psyche & NASA & Exploration d'un astéroïde métallique \\
Starlink & SpaceX & Maintien d'orbite des constellations \\
\bottomrule
\end{tabular}
\end{table}

\begin{figure}[H]
\centering
\includegraphics[width=0.8\textwidth]{images/bepi-colombo.png}
\caption{La mission BepiColombo (ESA/JAXA) utilise quatre propulseurs à effet Hall de type T6 pour son voyage de 7 ans vers Mercure, démontrant la fiabilité de cette technologie pour les missions longue durée.}
\label{fig:bepi-colombo}
\end{figure}

\subsection{Avantages par rapport à la propulsion chimique}

\begin{itemize}
    \item \textbf{$I_{sp}$ élevée} : 5 à 10 fois supérieure, donc bien moins de carburant nécessaire.
    \item \textbf{Durée de vie opérationnelle} : fonctionnement continu sur des milliers d'heures pour les missions longues.
    \item \textbf{Précision} : poussée finement contrôlable pour des manœuvres orbitales précises.
    \item \textbf{Réduction de masse} : permet d'embarquer plus de charge utile ou d'allonger la durée de mission.
\end{itemize}

%% ============================================================
%% VII. LIMITES, VERROUS TECHNOLOGIQUES ET PERSPECTIVES
%% ============================================================
\section{Limites, verrous technologiques et perspectives}

\subsection{Défis actuels}

\begin{table}[H]
\centering
\caption{Principaux défis technologiques}
\begin{tabular}{>{\bfseries}l p{9cm}}
\toprule
Défi & Description \\
\midrule
Érosion des parois & L'interaction du plasma avec les parois en céramique du canal provoque une érosion progressive (pulvérisation ionique) qui limite la durée de vie à 5\,000--10\,000~h. \\
Instabilités plasma & Les oscillations de décharge (modes basse et haute fréquence) peuvent dégrader les performances et compliquer la régulation. \\
Montée en puissance & Les propulseurs actuels fonctionnent jusqu'à 5--10~kW. L'objectif est d'atteindre 50--100~kW pour les missions cargo et l'exploration humaine. \\
Fiabilité électronique & Les composants de la PPU doivent résister aux radiations spatiales sur des durées de 15--20 ans. \\
\bottomrule
\end{tabular}
\end{table}

\subsection{Perspectives de développement}

\begin{itemize}
    \item \textbf{Propulsion haute puissance} : développement de propulseurs de 50--100~kW pour les missions de transport cargo vers Mars et les stations spatiales en orbite lunaire (Lunar Gateway).
    \item \textbf{Missions \textit{deep space}} : optimisation des trajectoires avec un $\Delta V$ cumulé élevé, rendant accessibles les planètes extérieures du système solaire.
    \item \textbf{Architectures hybrides} : combinaison propulsion chimique (forte poussée pour les manœuvres rapides) et électrique (haute $I_{sp}$ pour les transferts longs).
    \item \textbf{Nouveaux matériaux} : parois en carbure de silicium ou revêtements anti-érosion pour allonger la durée de vie.
    \item \textbf{Propergols alternatifs} : recherche sur le krypton (moins cher que le xénon) et l'iode (stockable sous forme solide).
\end{itemize}

%% ============================================================
%% CONCLUSION
%% ============================================================
\section{Conclusion}

Les moteurs à effet Hall illustrent la convergence entre physique des plasmas et électronique de puissance. Leur principe repose sur un mécanisme élégant : le piégeage magnétique des électrons crée une zone d'ionisation efficace, tandis que le champ électrique axial accélère les ions à haute vitesse, produisant une poussée avec une efficacité en propergol très supérieure à la propulsion chimique.

Cependant, l'électronique de puissance ne constitue pas seulement un support fonctionnel : elle est un \textbf{élément structurant des performances globales} du propulseur. La PPU doit convertir, réguler et distribuer l'énergie avec précision, tout en gérant le comportement dynamique et non linéaire du plasma. Les avancées dans ce domaine -- rendement des convertisseurs, robustesse aux radiations, stratégies de contrôle -- conditionnent directement l'avenir de la propulsion électrique spatiale.

%% ============================================================
%% RÉFÉRENCES
%% ============================================================
\section{Références}

\subsection*{Ouvrages de référence}
\begin{enumerate}
    \item Goebel, D. M., Katz, I. (2008) -- \textit{Fundamentals of Electric Propulsion: Ion and Hall Thrusters}, Wiley.
    \item Jahn, R. G. (1968) -- \textit{Physics of Electric Propulsion}, McGraw-Hill.
\end{enumerate}

\subsection*{Articles -- Techniques de l'Ingénieur (accès CNAM)}
\begin{enumerate}
    \setcounter{enumi}{2}
    \item Mazouffre, S. (2018) -- \textit{Propulsion électrique pour les systèmes spatiaux}, Réf. TRP4051, DOI: 10.51257/a-v1-trp4051.
    \item Pace, L., Daire, B., Beley, M., Salomez, F. (2024) -- \textit{Conversion DC-DC aux fréquences VHF -- Enjeux, avancement et perspectives}, Réf. E3979, DOI: 10.51257/a-v1-e3979.
    \item Costa, F. (2010) -- \textit{Compatibilité électromagnétique CEM -- Présentation générale}, Réf. D1300, DOI: 10.51257/a-v2-d1300.
\end{enumerate}

\subsection*{Revues scientifiques}
\begin{itemize}
    \item IEEE Transactions on Plasma Science
    \item IEEE Transactions on Power Electronics
    \item Journal of Propulsion and Power (AIAA)
\end{itemize}

\subsection*{Ressources en ligne}
\begin{itemize}
    \item NASA Technical Reports Server : \url{https://ntrs.nasa.gov/}
    \item ESA -- Electric Propulsion : \url{https://www.esa.int/}
\end{itemize}

\end{document}
